\documentclass[pdflatex,sn-nature]{sn-jnl}

\usepackage{graphicx}
\usepackage{amsmath,amssymb}
\usepackage{booktabs}
\usepackage[T1]{fontenc}
\usepackage{lmodern}
\usepackage{textcomp}
\usepackage{pdflscape}

\raggedbottom
\makeatletter
\newcommand\tabinput[1]{\@@input #1 }
\makeatother
\renewcommand{\tablename}{Supplementary Table}
\renewcommand{\thetable}{\arabic{table}}

\begin{document}

\title[Supplementary Information]{Supplementary Information for: Shared patterns of decoder drift and electrode failure
suggest simple safeguards for implanted brain--computer interfaces}

\author*[1]{\fnm{Md. Imtiaj Alam} \sur{Sajin}}\email{imtiajsajin@gmail.com}
\author[1]{\fnm{Esm E Moula} \sur{Chowdhury Abha}}\email{esmechowdhuryabha@gmail.com}
\author[1]{\fnm{Md Wahiduzzaman} \sur{Suva}}\email{wshuvo360@gmail.com}
\affil*[1]{\orgdiv{Department of Computer Science}, \orgname{American International University-Bangladesh},
\orgaddress{\city{Dhaka}, \postcode{1229}, \country{Bangladesh}}}

\abstract{This document contains Supplementary Notes 1--3 and Supplementary Tables 1--12. Every table is generated
from the released result files by \texttt{scripts/make\_supp\_tables.py}, so the values match the code repository
exactly.}

\maketitle

\section*{Supplementary Note 1: Label-free monitoring statistics}

Before the main analyses, we asked whether label-free statistics computed on a later day can predict how well a
renormalized fixed decoder will perform on that day. The pilot used the first 46--61 sessions of monkey N
(January 2020 to about March 2021) and 283 session pairs at most 60 days apart. Decoders were ridge regressions on 8
lags of spike-band power ($\alpha = 0.1$), trained on the first 300 trials and tested on the rest.

We computed 14 statistics that need no labels on the later day:
\begin{itemize}
\item the shift in channel means and in channel standard deviations;
\item the principal angles between the dominant subspaces of the two days;
\item disagreement within a bootstrap ensemble of decoders;
\item disagreement between decoders built on different views of the data (spike-band power versus threshold
crossings, and 8 lags versus 1 lag);
\item statistics of the predicted output (velocity bias, variance ratio, Wasserstein distance to the training
kinematics and smoothness);
\item self-consistency between predicted position and predicted velocity;
\item Procrustes alignment disagreement, the Procrustes residual and the share of decoder readout outside the
dominant subspace;
\item neural modulation depth;
\item the log ratio of impedances and the number of silent channels.
\end{itemize}

Elapsed time predicted decoder accuracy well (Spearman $\rho = -0.75$ across all pairs and $-0.56$ within 60 days).
After controlling for log elapsed days, the partial Spearman correlation of every statistic with $R^2$ was at most
0.27 (output variance) and mostly below 0.1. In time-blocked cross-validated prediction of $R^2$, elapsed days alone
gave a mean absolute error of 0.064. The label-free statistics alone gave 0.067--0.071, and elapsed days together with
the statistics gave 0.067--0.070. When the target was the gain from recalibration rather than $R^2$, elapsed days gave
0.045 and the statistics gave 0.054. For detecting unexpectedly poor days (residual more than one standard deviation
below the time trend), every statistic had an area under the receiver operating characteristic curve of 0.50--0.58.

The residual variation was not noise. Split-half reliability of the residuals, computed on alternating pairs of
trials, was $r = 0.72$, or 0.84 after Spearman--Brown correction. Alternating pairs were needed because centre-out
trials alternate between outward and return targets, which made odd/even single-trial splits negatively correlated.
Adding the reference decoder's own accuracy on the later day, a label-based measure of day difficulty, raised the
explained variance from 0.29 (elapsed days alone) to 0.56. Day-to-day performance beyond the time trend is therefore
real and reliable, but none of the label-free statistics we tried captured it. Abrupt faults on single channels
remain detectable by construction, and such faults may still justify monitoring.

\section*{Supplementary Note 2: Simulator}

The simulator (\texttt{ibci.sim}) is a feature-level generative model of chronic change. It transforms a real base
session into a simulated session $\Delta t$ days later, so that behaviour and latent dynamics stay realistic. Its
components are listed in Supplementary Table~11 and defined in the Methods. Two further components, heavy-tailed gain
jumps and random walks of raw offsets and scales, are included for decoders that are not renormalized. They do not
affect renormalized decoders and were not calibrated.

\textit{Calibration.} The simulated ladder must be scored with exactly the same correction code and regularization as
the real ladder. An earlier calibration used regularization 100--1,000 times stronger than the values chosen on real
data. That handicapped the simulated corrections, and the optimizer compensated with distorted parameters (loss 0.663
versus 0.074 after the fix). We also restricted silencing to channels that were active in the base session, because
most channels are already silent on any given day.

\textit{Validation.} The simulator calibrated on days before 700 was evaluated on base sessions after day 700. The
comparison covered all rungs of the ladder and labelled-data budgets of 20, 50 and 100 trials. A simulator with all
drift switched off served as a reference. The calibrated simulator was closer to the real data in every group (median
absolute error 0.12, 0.06 and 0.11 for fitted rungs, unfitted rungs and unfitted budgets; no-drift reference 0.37,
0.27 and 0.39). Its systematic miss, faster decay than observed in the later years, led to the separate late-period
calibration in Supplementary Table~11.

\textit{Intended uses and limits.} The simulator is meant for stress-testing decoders and label-free stabilizers under
controlled drift and electrode failure, and for comparing recalibration policies. It is calibrated on one monkey with
Utah arrays, spike-band power and a finger task. Its mechanisms reproduce the measured statistics, which does not show
that they are the true physical mechanisms. Users can calibrate it to other datasets with
\texttt{scripts/calibrate\_sim.py}.

\section*{Supplementary Note 3: Analysis choices that were corrected during the study}

We report three analysis choices that we revised, because each could mislead similar studies.

\textit{A full input remap is not a structural test.} A free $C\times C$ linear map in front of a frozen decoder
appears to test whether drift is a ``remixing'' of channels. With a readout of rank at most 32 (8 lags by 4 outputs)
and 96 channels, however, the map can reproduce any decoder of the same form. Its recovery therefore matches
retraining and says nothing about the structure of drift. We used it only as a calibration target for the simulator.

\textit{Recovery from re-learning a few channels reflects decoder importance.} Re-learning the weights of $k$ channels
recovered much of the loss when the channels were those the old decoder relied on most. Channels chosen by label-free
change did worse (Supplementary Table~9). The recovery therefore reflects which channels matter to the decoder, not
evidence that drift is concentrated on a few electrodes.

\textit{Small-budget recalibration must shrink the intercept.} Refitting the intercept freely with 10--20 trials
shifted the output offset and made recalibration harmful in many pairs. Shrinking the intercept toward the old
decoder, together with the weights, removed most of this harm (Supplementary Table~8).

\clearpage

\begin{landscape}
\begin{table}[p]
\caption{\textbf{Datasets.} Session counts are per analysis. Some human sessions contribute to more than one
analysis.}
\label{tab:data}
\scriptsize\setlength{\tabcolsep}{3.5pt}
\begin{tabular}{@{}llllll@{}}
\toprule
Individual & Source & Arrays and site & Signal & Task & Sessions (span) \\
\midrule
Monkey N & DANDI 001201 (LINK) & 2 Utah, motor cortex & SBP, TC & finger flexion & 312 on 303 days (3.4 years) \\
Monkey C & DANDI 000688 & Utah, motor/premotor & sorted spikes & reaching & 68 (3.0 years) \\
Monkey M & DANDI 000688 & Utah, motor/premotor & sorted spikes & reaching & 28 (1.5 years) \\
Human T6 & Dryad x0k6djj1h & 1 Utah, motor cortex & SBP & closed-loop cursor & 124 (3.1 years) \\
14 humans & Dryad x0k6djj1h & 20 Utah arrays & TC rates, impedance & research sessions & 2,289 ($\leq$7.3 years) \\
T11, T5 & Dryad n2z34tn5s (MINDFUL) & motor cortex & SBP (T11), TC (T5) & closed-loop cursor & 15 and 6 days \\
\botrule
\end{tabular}
\footnotetext{SBP, spike-band power; TC, threshold crossings.}
\end{table}
\end{landscape}

\begin{landscape}
\begin{table}[p]
\caption{\textbf{Correction ladder by individual and gap.} Median decoding $R^2$ for each rung, and median retention
of the renormalized fixed decoder (renormalized $R^2$ divided by reference $R^2$) with 95\% cluster-bootstrap
confidence intervals. Each training session contributes at most one pair per gap. Supervised rungs used 300
labelled trials.}
\label{tab:ladder}
\scriptsize\setlength{\tabcolsep}{3.5pt}
\begin{tabular}{@{}rrrrrrrrrrl@{}}
\toprule
Gap (days) & Pairs & Reference & None & Mean update & Renorm. & Aligned & Stable aligned & Gains & Ridge to prior &
Retention [95\% CI] \\
\midrule
\tabinput{supp/tab_ladder.tex}
\botrule
\end{tabular}
\end{table}
\end{landscape}

\begin{table}[p]
\caption{\textbf{Recurrent versus linear decoder (monkey N, 79 matched pairs).} Median reference $R^2$ (decoder
trained on the test day), retention of the renormalized fixed decoder and retention after re-learning channel gains.}
\label{tab:lstm}
\scriptsize\setlength{\tabcolsep}{3.5pt}
\begin{tabular}{@{}rrrrrrrr@{}}
\toprule
& & \multicolumn{2}{c}{Reference $R^2$} & \multicolumn{2}{c}{Retention, renormalized} &
\multicolumn{2}{c}{Retention, gains} \\
\cmidrule{3-4}\cmidrule{5-6}\cmidrule{7-8}
Gap (days) & Pairs & Ridge & LSTM & Ridge & LSTM & Ridge & LSTM \\
\midrule
\tabinput{supp/tab_lstm.tex}
\botrule
\end{tabular}
\end{table}

\begin{table}[p]
\caption{\textbf{Offline monitoring, monkey N (537 pairs).} Spearman correlation with decoder $R^2$, raw and partial
given log elapsed days, and time-blocked cross-validated error of predicted $R^2$. ``Days and divergence'' uses log
days together with the neural and output divergences.}
\label{tab:mindful_link}
\scriptsize\setlength{\tabcolsep}{3.5pt}
\begin{tabular}{@{}lrr@{}}
\toprule
Predictor & Raw $\rho$ & Partial $\rho$ \\
\midrule
\tabinput{supp/tab_mindful_link.tex}
\botrule
\end{tabular}
\end{table}

\begin{landscape}
\begin{table}[p]
\caption{\textbf{Closed-loop monitoring in the public MINDFUL data.} Pearson correlation of angular error with elapsed
days and with each divergence (raw and partial given days), and leave-one-day-out mean absolute error (degrees) of
predicted angular error from elapsed days alone, with neural or output divergence added, or from output divergence
alone.}
\label{tab:mindful}
\scriptsize\setlength{\tabcolsep}{3.5pt}
\begin{tabular}{@{}lrrrrrrrrrrr@{}}
\toprule
& & & & \multicolumn{2}{c}{Neural divergence} & \multicolumn{2}{c}{Output divergence} &
\multicolumn{4}{c}{Prediction error (\textdegree)} \\
\cmidrule{5-6}\cmidrule{7-8}\cmidrule{9-12}
Participant & Days & Windows & $r$(error, days) & Raw & Partial & Raw & Partial & Days & + neural & + output &
Output only \\
\midrule
\tabinput{supp/tab_mindful.tex}
\botrule
\end{tabular}
\end{table}
\end{landscape}

\begin{landscape}
\begin{table}[p]
\caption{\textbf{Regularization by the 1\% rule.} The rule chose the largest ridge penalty whose median same-day $R^2$
was within 1\% of the best. Columns give median $R^2$ with $\alpha = 0.1$ and with the chosen $\alpha$, the median
paired gain with 95\% cluster-bootstrap confidence interval, and the share of training sessions whose mean gain was
positive.}
\label{tab:reg}
\scriptsize\setlength{\tabcolsep}{3.5pt}
\begin{tabular}{@{}llrrrrlr@{}}
\toprule
Individual & Chosen $\alpha$ & Gap (days) & Pairs & $R^2$, $\alpha = 0.1$ & $R^2$, rule & Gain [95\% CI] &
Share positive \\
\midrule
\tabinput{supp/tab_reg.tex}
\botrule
\end{tabular}
\end{table}
\end{landscape}

\begin{landscape}
\begin{table}[p]
\caption{\textbf{Recalibration with few labelled trials (monkey N, 86 pairs).} Median retention of the renormalized
fixed decoder (no labels), and for each number of labelled trials the median retention of ridge to prior / retraining
from scratch / channel gains. Regularization chosen by cross-validation on the labelled trials.}
\label{tab:eff}
\scriptsize\setlength{\tabcolsep}{3.5pt}
\begin{tabular}{@{}rrrlllll@{}}
\toprule
Gap (days) & Pairs & No labels & $n = 10$ & $n = 20$ & $n = 50$ & $n = 100$ & $n = 300$ \\
\midrule
\tabinput{supp/tab_efficiency.tex}
\botrule
\end{tabular}
\end{table}
\end{landscape}

\begin{table}[p]
\caption{\textbf{Choosing the shrinkage strength.} Share of harmful recalibrations (more than 0.01 $R^2$ below the
renormalized fixed decoder) when the strength of ridge to prior is chosen by cross-validation on the available trials
or from past sessions, with two-sided Fisher's exact test, and median retention for both policies and for the best
strength in hindsight (an upper bound, not a usable policy).}
\label{tab:policy}
\scriptsize\setlength{\tabcolsep}{3.5pt}
\begin{tabular}{@{}rrrrrrrr@{}}
\toprule
& & \multicolumn{3}{c}{Harmful (\%)} & \multicolumn{3}{c}{Median retention} \\
\cmidrule{3-5}\cmidrule{6-8}
Trials & Pairs & Cross-validation & History & $P$ & Cross-validation & History & Hindsight \\
\midrule
\tabinput{supp/tab_policy.tex}
\botrule
\end{tabular}
\end{table}

\begin{landscape}
\begin{table}[p]
\caption{\textbf{Re-learning 32 channels only.} Median retention of a full refit (ridge to prior, strength from past
sessions) and median paired difference in $R^2$ from it, with 95\% cluster-bootstrap confidence intervals, when only
32 channels are re-learned. Channels were chosen by decoder importance, by label-free change, by importance-weighted
change, or at random.}
\label{tab:targeted}
\scriptsize\setlength{\tabcolsep}{3.5pt}
\begin{tabular}{@{}rrrllll@{}}
\toprule
Trials & Pairs & Full refit & Importance & Label-free change & Weighted change & Random \\
\midrule
\tabinput{supp/tab_targeted.tex}
\botrule
\end{tabular}
\end{table}
\end{landscape}

\begin{landscape}
\begin{table}[p]
\caption{\textbf{Electrode failure statistics per array.} Part., participant; Sess., sessions. Yield is the
percentage of electrodes with a threshold-crossing rate of at least 2~Hz, as the median of the first / last three
sessions. $h_{\mathrm{off}}$ and $h_{\mathrm{on}}$ are the active-to-silent and silent-to-active switching rates per
1,000 days. Rev./sil. counts revived and silenced initially active electrodes. Kurt., excess kurtosis. Kurtosis is computed on session-to-session
changes in log rate after removing session-wide changes. Imp. $\rho$ is the Spearman correlation of the array
median with day. Edge $-$ int. gives the median edge minus median interior slope of log rate per year, and $P$ is
the one-sided Mann--Whitney test that edge electrodes decline faster. For monkey N, values are given for the lateral
and medial arrays where they differ.}
\label{tab:failure}
\scriptsize\setlength{\tabcolsep}{3.5pt}
\begin{tabular}{@{}llrrrrrrrrrr@{}}
\toprule
Part. & Array & Sess. & Span (d) & Yield (\%) & $h_{\mathrm{off}}$ & $h_{\mathrm{on}}$ & Rev./sil. & Kurt. &
Imp. $\rho$ & Edge $-$ int. & $P$ \\
\midrule
\tabinput{supp/tab_failure.tex}
\botrule
\end{tabular}
\end{table}
\end{landscape}

\begin{table}[p]
\caption{\textbf{Simulator parameters.} Values calibrated on monkey N sessions before day 700 (early) and from day 700
onwards (late). Switching rates were measured from electrode activity in each period rather than calibrated. Fixed
parameters: local mixing share $\beta = 0.5$, turnover concentration 5, 20 latent dimensions and silent-signal fraction
$\kappa = 0.37$.}
\label{tab:sim}
\scriptsize\setlength{\tabcolsep}{3.5pt}
\begin{tabular}{@{}llrr@{}}
\toprule
Parameter & Unit & Early & Late \\
\midrule
\tabinput{supp/tab_sim.tex}
\botrule
\end{tabular}
\end{table}

\begin{landscape}
\begin{table}[p]
\caption{\textbf{Training on simulated futures.} Median $R^2$ on real later sessions of monkey N for ridge decoders
trained on sessions after day 700 (simulator calibrated before day 700). Base, $\alpha = 0.1$; stronger ridge,
$\alpha = 10^4$ chosen on the calibration period; perturbation, random channel dropout and gain noise; simulator,
eight simulated futures; simulator and ridge, both; several days, real data from the three previous sessions;
reference, decoder trained on the test day.}
\label{tab:augment}
\scriptsize\setlength{\tabcolsep}{3.5pt}
\begin{tabular}{@{}rrrrrrrrr@{}}
\toprule
Gap (days) & Pairs & Base & Stronger ridge & Perturbation & Simulator & Simulator and ridge & Several days &
Reference \\
\midrule
\tabinput{supp/tab_augment.tex}
\botrule
\end{tabular}
\end{table}
\end{landscape}

\end{document}
